\documentclass[12pt]{article}
\usepackage{amsmath, amssymb, geometry, enumitem}
\geometry{margin=1in}
\setlength{\parindent}{0pt}
\setlength{\parskip}{0.5em}

\title{ECON 101: Macroeconomics \\ Solutions -- Quiz 5}
\author{Instructor: Muhebullah Karimzada}
\date{April 07,  2026}

\begin{document}
\maketitle

\begin{enumerate}[label=\textbf{\arabic*.}]

\item \textbf{Correct answer: (C)}\\
The Bank of Canada's main monetary policy tool is the \textbf{target overnight interest rate}. This is the interest rate it influences directly in order to affect other market interest rates, borrowing, spending, and inflation.

\item \textbf{Correct answer: (C)}\\
A decrease in the target overnight rate lowers short-term interest rates in the economy. Lower interest rates reduce borrowing costs for households and firms, which stimulates consumption and investment spending.

\item \textbf{Correct answer: (B)}\\
The monetary transmission mechanism works partly through interest rates. When interest rates change, firms adjust investment and households adjust spending, which then affects aggregate demand.

\item \textbf{Correct answer: (B)}\\
If Canadian interest rates rise, Canadian financial assets become more attractive to foreign investors. Demand for the Canadian dollar rises, so the dollar tends to \textbf{appreciate}.

\item \textbf{Correct answer: (C)}\\
Contractionary monetary policy raises interest rates, which reduces consumption, investment, and net exports. This lowers aggregate demand, so the AD curve shifts to the \textbf{left}.

\end{enumerate}


\textbf{Given:}
\begin{itemize}
    \item Potential GDP: \(Y^* = 2{,}500\)
    \item Initial actual GDP: \(Y = 2{,}700\)
    \item Investment falls by \(\$100\) billion
    \item MPC \(= 0.75\)
\end{itemize}

\subsection*{(a) Calculate the spending multiplier (3 marks)}

The simple spending multiplier is:

\[
k = \frac{1}{1-\text{MPC}}
\]

Substitute the given MPC:

\[
k = \frac{1}{1-0.75}
\]

\[
k = \frac{1}{0.25} = 4
\]

\[
\boxed{k = 4}
\]

\textbf{Answer:} The spending multiplier is \(\boxed{4}\).

\subsection*{(b) Calculate the total change in equilibrium GDP (4 marks)}

The formula for the total change in equilibrium GDP is:

\[
\Delta Y = k \times \Delta I
\]

We know:
\[
k = 4, \qquad \Delta I = -100
\]

Substitute:

\[
\Delta Y = 4 \times (-100)
\]

\[
\Delta Y = -400
\]

\[
\boxed{\Delta Y = -400}
\]

\textbf{Answer:} Equilibrium GDP decreases by \(\boxed{\$400 \text{ billion}}\).

\subsection*{(c) Calculate the new level of equilibrium GDP (2 marks)}

Initial GDP:
\[
Y_0 = 2{,}700
\]

Change in GDP:
\[
\Delta Y = -400
\]

New equilibrium GDP:

\[
Y_1 = Y_0 + \Delta Y
\]

\[
Y_1 = 2{,}700 - 400 = 2{,}300
\]

\[
\boxed{Y_1 = 2{,}300}
\]

\textbf{Answer:} The new equilibrium GDP is \(\boxed{2{,}300}\).

\subsection*{(d) AD--AS diagram explanation (3 marks)}

Students should draw an AD--AS diagram with:

\begin{itemize}
    \item a downward-sloping \(\text{AD}_0\) curve,
    \item an upward-sloping SRAS curve,
    \item a vertical LRAS curve at \(Y^* = 2{,}500\),
    \item an initial equilibrium where \(\text{AD}_0\) intersects SRAS at \(Y = 2{,}700\),
    \item a leftward shift of AD from \(\text{AD}_0\) to \(\text{AD}_1\),
    \item a new short-run equilibrium at \(Y = 2{,}300\).
\end{itemize}

\textbf{Interpretation:}
\begin{itemize}
    \item Initially, the economy is producing above potential GDP, so there is an \textbf{inflationary gap}.
    \item The increase in the policy rate reduces investment, which shifts AD left.
    \item In the short run, both output and the price level fall.
\end{itemize}
\newpage
\textbf{Required labels:}
\begin{itemize}
    \item AD\(_0\), AD\(_1\)
    \item SRAS
    \item LRAS
    \item Initial equilibrium output \(2{,}700\)
    \item New equilibrium output \(2{,}300\)
\end{itemize}

\subsection*{(e) Explain how this monetary policy reduces inflation in the medium run (3 marks)}

Step-by-step:

\begin{enumerate}
    \item The Bank of Canada raises the policy interest rate.
    \item Higher interest rates reduce borrowing and lower investment spending.
    \item Lower investment reduces aggregate demand.
    \item The AD curve shifts left.
    \item Output falls from \(2{,}700\) to \(2{,}300\).
    \item Because the economy was initially above potential GDP (\(2{,}700 > 2{,}500\)), there was upward pressure on prices.
    \item Lower aggregate demand reduces this inflationary pressure.
    \item In the medium run, slower demand growth helps inflation move back toward the Bank of Canada's target.
\end{enumerate}

\textbf{Important point:}  
This policy helps reduce inflation because it closes the inflationary gap and weakens excess demand in the economy.

\[
\boxed{\text{Contractionary monetary policy reduces inflation by lowering aggregate demand.}}
\]

\hrule
\vspace{1em}

\begin{center}
\textit{End of Solution Key}
\end{center}

\end{document}